\begin{lemma}[Nonempty support]
Every $X\in\mathsf{PowerQ}(Q)$ has nonempty support:
\[
\mathsf{PowerSupport}(X)\ne\varnothing.
\]
\end{lemma}
\begin{proof}
Proceed by induction on the inductive presentation from
\noderef{PowerQ}.  If $X=\mathsf{atom}(q)$, then the support equation in
\noderef{PowerSupportEquations} identifies its support with $\{q\}$, which
contains $q$.  If $X=\mathsf{node}(\iota,h,f)$, choose an index
$i_0\in\iota$ using the supplied witness $h$.  The induction hypothesis
gives an atom $q$ in the support of $f(i_0)$.  The node support equation in
\noderef{PowerSupportEquations} says that this support is contained in the
indexed union defining the support of $X$, so the same $q$ belongs to the
support of $X$.  Thus every atom and every nonempty node has nonempty
support.
\end{proof}
